\section{Odd primes and the ramification index}
\label{sec:odd-ramification}

A finite coefficient algebra makes the cyclotomic distance odd in a
precise range of ramification indices.  Keeping the entire algebra,
and selecting one of its residue factors by an endomorphism, makes
the argument work without a completeness assumption.

\begin{theorem}\label{thm:odd-ramification}
Let $p$ be an odd prime and let $O$ be any characteristic-zero discrete
valuation ring with residue field $k$ of characteristic $p$.
Write $e=v_O(p)$ for its absolute ramification index, and let $v_2$
denote the exponent of two in a positive integer.  Suppose
\[
 v_2(e)\le v_2(p-1).
\]
For $A=O[\epsilon]/(\epsilon^2)$, the kernel of
\[
 K_0(\Lat_A(S_p))\longrightarrow K_0(\Lat_O(S_p))
\]
contains a split countable direct sum of copies of $\Z/2$.
The displayed classes also vanish after inverting $p$.
In particular $\Phi_A$ is not surjective.
There is no completeness, finiteness, or perfection hypothesis on $k$.
\end{theorem}

\begin{corollary}\label{thm:even-ramified-cubic}
Let $\Z_3$ denote the complete ring of $3$-adic integers.
For $O=\Z_3[\zeta_3]$ or $\Z_3[\sqrt3]$, the flat dual numbers
$O[\epsilon]/(\epsilon^2)$ have a split countable degree-three
reduction-kernel subgroup isomorphic to $\bigoplus_{\mathbb N}\Z/2$.
\end{corollary}

\begin{proof}[Proof of the corollary]
Both rings have absolute ramification index two, so the theorem applies.
\end{proof}

\begin{proof}[Proof of Theorem~\ref{thm:odd-ramification}]
\emph{Primary composition lengths.}
For an $A$-lattice $E$ restricted to $C_p$, the two Tate groups of
$E_0=E/\epsilon E$ have finite $O$-length and are killed by $p$.
The canonical coefficient boundary $\delta$ is natural and
anticommutes with the connecting maps of admissible exact sequences,
by the coefficient calculation in
Proposition~\ref{prop:dual-bockstein-invariants}.
Thus $u=\delta^2$ commutes with the entire Tate exact sequence.
For a monic irreducible $f\in k[U]$, $f\ne U$, take the primary part
for the maximal ideal $(\pi,f)$ of $O[U]$, where $U$ acts by $u$.
Primary extraction is exact on finite-length modules: apply the
Chinese remainder decomposition for powers of their finitely many
supporting maximal ideals.
On this part $u$ is invertible, so $\delta$ identifies both parities
and the two connecting-image modules.  The six-term length formula
therefore proves that
\[
 \nu_f(E)=\operatorname{length}_{O[U]}H^0(E_0)_f\pmod2
\]
is additive.  The length is composition length over $O[U]$, not
dimension over $k$ or length over $O$.
Every constant coefficient extension has $\delta=0$.
The comparison~\eqref{eq:dual-prime-scaling} shows that changing the
generator by $g\mapsto g^a$ scales the residual primary parameter
by $a\in\F_p^\times$.

\emph{A finite coefficient algebra.}
Put $d=\gcd(e,p-1)$, $t=(p-1)/d$, and $m=e/d$.
The hypothesis makes $m$ odd.
Choose a uniformizer $\pi$ and form
\[
 O_1=O[s]/(s^t-\pi).
\]
The polynomial is Eisenstein, so this is a finite free domain over
$O$ of rank $t$.  Its unique maximal ideal is $(s)$; hence it is a
discrete valuation ring, with $v_s(p)=(p-1)m$.
Write $b=p/s^{(p-1)m}$, a unit, and define
\[
 G(Y)=b+\sum_{j=1}^{p-2}\frac{\binom p{j+1}}p\,
              b s^{mj}Y^j+Y^{p-1},\qquad C=O_1[Y]/G(Y).
\]
This is a finite free coefficient algebra over $O$.
Its residue polynomial $\bar b+Y^{p-1}$ is separable with nonzero
roots, so $C/sC$ is a product of finite separable fields over $k$.
The element $Y$ is a unit, and $s$ is a nonzerodivisor.
The defining polynomial gives literal identities
\[
 \zeta=1+s^mY,\qquad \Phi_p(\zeta)=0,\qquad
 \zeta^p=1,\qquad \zeta-1=s^mY.
\]
The algebra $C$ may be semilocal.  No selection of a local coefficient
factor, and no Hensel lifting, is required.

\emph{A companion lattice and its two boundaries.}
The following construction applies to this $C$ and to the analogous
algebras after further finite unramified coefficient extensions.
Write
\[
 h(Z)=\frac{Z^p-1}{Z-\zeta}=(Z-1)q(Z),\qquad
 N(Z)=1+Z+\cdots+Z^{p-1}=(Z-\zeta)q(Z).
\]
These are polynomial identities.  For example $h(1)=0$ follows from
$(1-\zeta)h(1)=0$ because $1-\zeta$ is a nonzerodivisor.
Let $V=C[Z]/h(Z)$, let $B$ be multiplication by $Z$ in its monomial
basis, let $v=1\in V$, and let $w$ extract the coefficient of
$Z^{p-2}$.  The polynomial $q$ is monic of degree $p-2$, so $w(q)=1$.
For any $c\in C$, define
\[
 T_c=\begin{pmatrix}B&\epsilon v\\\epsilon c w&\zeta\end{pmatrix}
\]
on a free $C[\epsilon]/(\epsilon^2)$-module of rank $p$.
Its first-order off-diagonal terms in $T_c^p$ are multiples of
$h(B)=0$; hence $T_c^p=I$.
There is no unit condition on $c$.

For the constant reduction $M=V\oplus C_\zeta$, the invariant
submodule of $V$ is $Cq$, and its norm image is $(1-\zeta)Cq$.
Indeed $N(B)f=(1-\zeta)f(1)q$.
Its norm kernel is $(B-1)V$, because $1-\zeta$ is a nonzerodivisor.
The scalar $\zeta$ summand has zero even Tate group and odd group
$C/(\zeta-1)$.  Consequently
\[
 H^0(M)=H^1(M)=C/(s^m).
\]
The even coefficient boundary sends $q$ to $cw(q)=c$.
The odd coefficient boundary sends $1$ to $q$, by the polynomial
divided-difference identity
\[
 \frac{N(B)-N(\zeta)}{B-\zeta}\,v=q(B)v=q.
\]
No inverse of $B-\zeta$ is taken in this expression.
Thus $u=\delta^2$ is exactly multiplication by $c$.

The reduced class is regular: there is a $C$-split exact sequence
\[
 0\longrightarrow C_\zeta\xrightarrow{1\mapsto h}
 C[C_p]\longrightarrow C[Z]/h(Z)\longrightarrow0.
\]
The kernel is freely generated by the monic polynomial $h$.
There is also an exact order-two relation.
For any dual-number lattice $E$, let $E^-$ be its coefficient twist
by $\epsilon\mapsto-\epsilon$.
The admissible sequence
\[
 0\longrightarrow E^-\longrightarrow
 (C[\epsilon]/(\epsilon^2))\otimes_C\operatorname{Res}_C E
 \longrightarrow E\longrightarrow0
\]
has injection $x\mapsto\epsilon\otimes x-1\otimes\epsilon x$
and evaluation surjection.  Its kernel is a free direct summand in
an underlying coefficient basis.
The underlying $C$-lattice is an extension of two copies of $E_0$.
For $E=E_c$, the matrix $\operatorname{diag}(I_{p-1},-1)$ identifies
$E^-$ with $E$.  It follows that
\begin{equation}\label{eq:odd-order-two}
 2\bigl([E_c]-[(C[\epsilon]/(\epsilon^2))[C_p]]\bigr)=0.
\end{equation}

\emph{Infinite residue fields.}
Suppose $k$ is infinite, and choose one field factor $F$ of $C/sC$.
There are infinitely many primitive elements $\alpha\in F$ such
that $\alpha$ and $a\alpha$ are not $k$-conjugate for any
$a\in\F_p^\times\setminus\{1\}$.
To prove this, work in a normal closure of $F/k$.
Each unwanted equation $\sigma(\alpha)=a\alpha$ defines a proper
$k$-linear subspace of $F$, except for $(\sigma,a)=(1,1)$:
if it held identically, evaluation at $1$ would give $a=1$ and then
$\sigma=1$.
Avoid these finitely many subspaces.  Over an infinite field their
complement has infinitely many points; the equations with $a=1$
also exclude every nonprimitive element.
Choose a countable sequence $\alpha_i$ with pairwise disjoint
scalar-conjugacy orbits, each of which is finite.

Lift $(\alpha_i,0,\ldots,0)\in C/sC$ to $c_i\in C$, putting
$\alpha_i$ in the selected factor.  Let $f_i$ be its minimal
polynomial over $k$.
The $s$-filtration of $C/(s^m)$ has $m$ factors $C/sC$.
In each layer exactly one factor contributes to the $f_i$-primary
part: it is the simple $O[U]$-module $F$, with $U=\alpha_i$.
Every other factor has $U=0$ and disappears.
The primary composition length is therefore the odd integer $m$.

\emph{Finite residue fields.}
Suppose $k=\F_Q$, and choose $r$ such that $\bar G$ splits over
$\F_{Q^r}$.
For strictly increasing positive multiples $n_i$ of $r$, lift a
monic irreducible polynomial of degree $n_i$ over $k$ to $O_1$.
Its quotient $S_i$ is finite free over $O_1$ and has unique maximal
ideal $(s)$ and residue field $k_i=\F_{Q^{n_i}}$; thus it is a
discrete valuation ring.
Set $C_i=S_i[Y]/G(Y)$.
Its reduction is a product of $p-1$ copies of $k_i$.
Choose a generator $\alpha_i$ of $k_i^\times$, lift
$(\alpha_i,0,\ldots,0)$ to $c_i\in C_i$, and let $f_i$ be its
minimal polynomial over $k$.
Again the selected primary length is $m$.
The scalar-conjugacy orbit is free: if
$\alpha_i^{Q^j}=a\alpha_i$ with $0<j<n_i$ and $a\in\F_p^\times$,
then
\[
 Q^{n_i}-1\mid(p-1)(Q^j-1),
\]
although the positive right side is smaller than the left side.
For $j=0$ one obtains $a=1$.
Distinct degrees make the chosen scalar orbits disjoint.

\emph{Induction and splitting.}
In either case write $C_i$ for the chosen coefficient algebra and
$a_i$ for its rank over $O$.
Restrict $E_{c_i}$ to $A$, induce from $C_p$ to $S_p$, and subtract
$a_i[A[S_p]]$; call the resulting class $x_i$.
Restriction preserves freeness and exactness because $C_i$ is finite
free over $O$.
The reduced regular-class equality proves that $x_i$ specializes to
zero over $O$, and~\eqref{eq:odd-order-two} gives $2x_i=0$.
In the Mackey restriction back to $C_p$, the normalizer terms give
each scalar twist $a\in\F_p^\times$ exactly once.
All other intersections are trivial and give Tate-acyclic free
group lattices.  Therefore
\[
 \bigl(\nu_{f_j}\circ\operatorname{Res}_{C_p}^{S_p}\bigr)(x_i)
 =\begin{cases}1,&i=j,\\0,&i\ne j.\end{cases}
\]
This proves nonzero exact order two and independence.
Every finite-length Tate module has finite primary support, so
the direct sum of these invariants retracts the inclusion of the
$x_i$.  The invariants vanish on every Young permutation class.
After inverting $p$, rational nil invariance identifies each class
with its regular reduction, proving the last vanishing assertion.
\end{proof}

\begin{remark}\label{rem:odd-ramification-scope}
The condition in Theorem~\ref{thm:odd-ramification} is sufficient.
When $v_2(e)>v_2(p-1)$, the cyclotomic distance $e/\gcd(e,p-1)$
in this construction is even, and its parity detectors vanish.
Further finite coefficient extensions cannot make that distance odd:
if $e'$ is a multiple of $e$ and $p-1$ divides $e'$, then
$e'/(p-1)$ remains even.
This does not establish a positive result in the complementary range.
\end{remark}
